\scanpage{84}
\uline{Calculs} pour $T_\ell(G_m)\inserted{=\mathbb{Q}_\ell(1)}$ et $G_m$.

$X$ de type fini sur $\mathbb{Z}$

\[
\begin{aligned}
H^0(X,\mathbb{Q}_\ell(1))\deleted{\unsure{\otimes G_m)}}&=0\\
H^1(X,\mathbb{Q}_\ell(1))&=\tikz[baseline=(a.base)]{\node[draw,ellipse,inner sep=1pt] (a) {$H^0(X,O_X^*)\otimes_{\mathbb{Z}}\mathbb{Q}_\ell$};}\\
H^2(X,\mathbb{Q}_\ell(1))&\quad\text{extension de }\,
\tikz[baseline=(b.base)]{\node[draw,rounded corners,inner sep=2pt] (b) {$T_\ell(H^2(X,G_m))$};\draw (b.north east) .. controls +(.13,.22) and +(-.18,-.04) .. ++(.55,.24);}
\quad\text{par}\\
&\hspace{4em}\tikz[baseline=(c.base)]{\node[ellipse,draw,inner sep=1pt] (c) {$\deleted{\unsure{R}}H^1(X,\unsure{\underline{O}_X^*|O_X^*|\mathcal{O}_X^*})\otimes\mathbb{Q}_\ell$};}
\end{aligned}
\]
\margin{\uline{coBrauer}}
$H^3(X,\mathbb{Q}_\ell(1))\simeq T_\ell(H^3(X,G_m))$ est ce qu'il peut [p. ex. $\mathbb{Z}_\ell$ si $X$ \deleted{\unsure{modèle}}modèle can. d'un corps de fonctions d'une variable sur $\mathbb{F}_p$]
\[
\begin{aligned}
H^i(X,\mathbb{Q}_\ell(1))&=0\qquad\text{si}\quad i\geq4\qquad ??\\
\shortparallel&\\
T_\ell(H^i(X,G_m))&
\end{aligned}
\]

On conjecture (Artin) que $H^2(X,G_m)={}\text{groupe fini}$ si $X$ propre sur $\mathbb{Z}$. \deleted{Mais si} \inserted{Vrai} si $X$ courbe projective lisse sur $\mathbb{F}_p$, mais faux quand on lui enlève des pts \textendash\ le rang du $H^2$ augmente d'autant

$H^3(X,\mathbb{Q}_\ell(1))$ : quel est sa taille [il décroît en passant à un ouvert de $X$]
% archive sheet number: 83
